\documentclass[a4papaer, 11pt]{article}
%패키지
\usepackage[scale=0.75,twoside,bindingoffset=5mm]{geometry}
\usepackage{kotex} %한글 사용
\usepackage{amsmath} %align 등 기능

\begin{document}
\title{Description of signal processing methods and descriptive statistics
\author{응용통계학과 석사과정 문종민\\
}
}


\maketitle
\subsubsection{Preprocessing of the EMG signal}

We used Matlab
(Mathworks, Natick, MA, USA).
to process the signal in offline.

\paragraph{High pass and notch filter}
To minimize the shift of the baseline,
3rd order Butterworth high pass filter was applied to the signal at 30Hz.
For the purpose of noise reduction,
we applied 5th order Butterworth notch filter at 60Hz.
For this step of processing,
we used the \texttt{butter} function in the Signal Processing Toolbox™ of Matlab.

\paragraph{Full-wave rectification}
Then we applied full-wave rectification to the signal, by taking absolute value of the signal data.

\paragraph{gernetation of envelopes}
We generated EMG envelopes of the rectified signal
by applying cubic spline interpolation over local maxima separated by at least \textit{[parameter used by the researcher e.g. 200, 350...]} samples. For this step we used \texttt{envelope} function with 'peak' option in the Signal Processing Toolbox™ of Matlab.


\subsubsection{Descriptive statistics}

\paragraph{Mean peak to peak}
This is average of difference between all local maxima and local minima of the signal. To find the local maxima and minima, we used \texttt{findpeaks} function which finds the local maxima of the signal. We multiplied $-1$ to the signal data and applied the \texttt{findpeaks} function to get the local maxima. Local maximas and minimas whose absolute value is too small were not used in the calculation, and the threshold value was chosen heuristically.
\paragraph{Mean RMS(or mean absolute value)}Mean RMS is the arithmetic mean of the absolute value of the signal within a specified range.

\paragraph{Signal to noise ratio(SNR)}
With pre-specified signal interval $A$ and noise interval $B$,
we calculate root mean square(RMS) of each interval.
Let $Psignal = RMS(A)$ and $Pnoise = RMS(B)$. Then we define as
\begin{align*}
	SNR = k * log_a(Psignal / Pnoise)
\end{align*}
Where k is chosen by the researcher and $a \in \{2, e, 10\}$.
To calculate the RMS, We used \texttt{rms} function in the Signal Processing Toolbox™ of Matlab.


\end{document}
